\chapter{Literature Survey}

The phase-ordering problem has been studied for more than four decades, first as a question of compiler engineering, then as a search problem, and most recently as a machine-learning problem. This chapter reviews the literature that shaped the present work. It is organized in six parts: the classical foundations that define the passes and first recognized their interactions; theoretical results on the nature of the problem; search-based approaches that treat ordering as an optimization problem in its own right; machine-learning approaches that try to predict good orderings; the infrastructure, benchmark suites, and evaluation methodology on which this project rests; and educational tools for making compiler internals visible. The chapter closes with a summary table and a statement of the gap that this project addresses.

\section*{Classical foundations}

Aho, Lam, Sethi, and Ullman \cite{aho2006} provide the standard treatment of the compiler pipeline that this project follows: lexical analysis, syntax analysis by recursive descent, semantic analysis with a scope-aware symbol table, generation of three-address code, and machine-independent optimization. Their chapters on code optimization define the passes implemented in this project, namely constant folding, copy propagation, dead code elimination, common subexpression elimination through available-expression analysis, and strength reduction, and present them as transformations over basic blocks and flow graphs. The book observes that ``the order in which optimizations are applied'' affects results and recommends iterating certain passes to a fixed point, but it does not attempt to characterize the interactions systematically. The three-address code used here, with at most one operator and two operands per instruction and explicit temporaries, follows their presentation closely, as does the treatment of short-circuit boolean expressions through jumping code.

Cooper and Torczon \cite{cooper2011} devote an explicit discussion to the sequencing of optimizations. They note that the compiler writer must decide both which transformations to include and the order in which to apply them, that some transformations create opportunities for others while some destroy them, and that the ideal sequence varies from program to program. They describe the fixed sequence as a compromise and point to their own research on adaptive compilation as an alternative. The book's treatment of local value numbering and of the distinction between local, superlocal, regional, and global optimization scope informed the decision in this project to keep every pass local to a basic block and to reset analysis state at labels and jumps, trading some optimization power for guaranteed safety under arbitrary reordering.

Muchnick \cite{muchnick1997} presents what is probably the best-known prescriptive answer to the phase-ordering question: a diagram that arranges dozens of optimizations into tiers and proposes an order in which they should be run, with constant folding and algebraic simplification early, common subexpression elimination and copy propagation in the middle, and dead code elimination repeatedly as a clean-up step. The order is justified by experience rather than by measurement, and Muchnick explicitly acknowledges that it is a heuristic. The empirical structure uncovered by this project, in which enabling passes must precede clean-up passes and dead code elimination is most effective when it runs last, agrees with Muchnick's recommendation and provides quantitative support for it on a small pass set.

Bacon, Graham, and Sharp \cite{bacon1994} catalogue the program transformations used by optimizing compilers for high-performance computing. Although the catalogue is oriented towards loop and data-layout transformations that are outside the scope of this project, it is valuable for two reasons. It defines each transformation precisely, including the scalar transformations implemented here, and it discusses the legality conditions that a transformation must respect, which guided the conservative handling of arrays and global variables in the passes: because no alias analysis is performed, array loads are never cached across instructions by common subexpression elimination, and array stores and global stores are never removed by dead code elimination.

\section*{The nature of pass interactions}

Whitfield and Soffa \cite{whitfield1997} were among the first to study pass interactions formally. They introduced a specification language, Gospel, in which the preconditions and effects of a code-improving transformation are written down explicitly, and a tool, Genesis, that generates the transformation from the specification. Using the specifications, they derived \emph{enabling} and \emph{disabling} relationships between pairs of transformations: an enabling relationship exists when applying one transformation can create the preconditions for another, and a disabling relationship exists when it can destroy them. Their experiments showed that these relationships are real and measurable, that many pairs interact in both directions, and that the interaction depends on the program. The pairwise interaction matrix produced by this project's visualizer, which compares the mean code size of orderings in which pass A precedes pass B with that of orderings in which B precedes A, is a direct empirical analogue of their enabling and disabling analysis.

Click and Cooper \cite{click1995} approached the interaction problem from a different angle. They showed that, for some pairs of analyses, running them separately in either order is strictly weaker than running a single combined analysis, because each analysis can only exploit the facts that the other has already discovered. Their combined constant propagation and unreachable-code elimination finds constants that no sequence of the two separate passes can find. This result is important for the interpretation of exhaustive phase-ordering studies such as the present one: it establishes that the best sequence of separate passes is not necessarily the best achievable result, and that some phase-ordering problems can only be dissolved by merging passes rather than reordering them. The present project deliberately kept the passes separate in order to study ordering, but this result is a natural direction for future work.

Touati and Barthou \cite{touati2006} examined the decidability of the phase-ordering problem. They modelled optimizations as functions over programs and asked whether an algorithm can decide, for an arbitrary program and performance goal, whether some sequence of optimizations reaches the goal. They proved that the problem is undecidable in the general model where sequences may be arbitrarily long and the performance function is arbitrary, and that it becomes decidable when the sequence length is bounded or when the set of reachable programs is finite. The result justifies the design choice made here: by fixing the set of passes at six and requiring each pass to run exactly once, the ordering space is finite and small, and exhaustive enumeration is not merely feasible but complete.

\section*{Search-based approaches}

Cooper, Schielke, and Subramanian \cite{cooper1999} were the first to treat pass ordering as an explicit optimization problem. Working in a research compiler with ten transformations and targeting code size for embedded systems, they used a genetic algorithm to evolve sequences of twelve passes and compared the result with the compiler's fixed default sequence. The evolved sequences produced substantially smaller code, and, more importantly, the best sequence differed from program to program. Their paper introduced the vocabulary of ``adaptive compilation'' and established code size, rather than execution speed, as a legitimate first-class objective, a framing that this project adopts through its emphasis on static instruction count for embedded targets.

Almagor, Cooper, Grosul, Harvey, Reeves, Subramanian, Torczon, and Waterman \cite{almagor2004} set out to characterize the shape of the search space rather than merely to search it. For a set of small programs they enumerated every sequence of ten passes drawn from five transformations, some $5^{10}$ sequences in all, and studied the distribution of outcomes. They found that the space contains many local minima, that good sequences are nevertheless common, and that simple search strategies such as hill climbing with random restarts find sequences within a few percent of the best in a fraction of the time an exhaustive search would take. The observation that the space, though large, is highly structured is the same observation that emerges from the present project at a smaller scale, where the 720 orderings of six passes collapse into at most 23 distinct outcomes for any single program.

Kulkarni, Hines, Hiser, Whalley, Davidson, and Jones \cite{kulkarni2004} attacked the cost of search. Working in the VISTA interactive compilation framework built on the VPO optimizer, they observed that genetic-algorithm searches over phase sequences wasted most of their time re-evaluating sequences whose outcomes were already known. Two techniques removed the waste: detecting \emph{dormant} phases, which find nothing to do and leave the code unchanged, so that any sequence differing only in dormant phases can be skipped; and hashing the resulting function bodies so that sequences that produce identical code are recognized and evaluated only once. Together these reduced search time by an order of magnitude. Both phenomena are directly visible in this project's data: strength reduction and algebraic simplification are dormant with respect to code size on most benchmarks, and the majority of the 720 orderings produce code identical to some other ordering.

Kulkarni, Whalley, Tyson, and Davidson \cite{kulkarni2006} then asked whether the space could be enumerated exhaustively rather than searched. Their insight was to enumerate not sequences but \emph{distinct function instances}: starting from the unoptimized function, they applied every phase, kept only the results that differed from anything already seen, and repeated until no new instance appeared. Because most sequences produce duplicates, the space of instances is orders of magnitude smaller than the space of sequences. With fifteen reorderable phases they were able to enumerate the complete space of instances for the large majority of functions in their embedded benchmark set, and to determine the truly optimal sequence for each. The exhaustive design of the present project, which enumerates every permutation of a smaller pass set, is a simplified realization of the same idea; the equivalence-class analysis in Chapter 5 is the instance-space view of the same data.

Kulkarni, Whalley, Tyson, and Davidson \cite{kulkarni2009} extended the exhaustive approach to dynamic performance. Because measuring the execution time of every function instance by simulation is expensive, they measured dynamic instruction counts by simulation only for instances with distinct control-flow structure and estimated the rest, achieving near-perfect correlation with full simulation at a small fraction of the cost. They were then able to report, for each function, how far the compiler's default sequence fell short of the optimal one on both size and speed. The present project measures dynamic instruction counts exactly by interpreting every ordering, which is affordable at its scale, and confirms their qualitative finding that size and speed are strongly but not perfectly correlated across orderings.

Triantafyllis, Vachharajani, Vachharajani, and August \cite{triantafyllis2003} proposed Optimization-Space Exploration, in which the compiler compiles each code segment under several optimization configurations and selects the best by a static performance estimate rather than by execution. The configurations are pruned in advance to a small set that covers most of the benefit, and the selection is done at compile time so that the approach fits inside a production compiler. Their use of a static estimator to rank configurations is the same strategy as this project's weighted cycle model; their finding that a small number of well-chosen configurations captures most of the available benefit parallels the observation here that a handful of ordering classes account for all the distinct outcomes.

Hoste and Eeckhout \cite{hoste2008} introduced COLE, Compiler Optimization Level Exploration, which uses a multi-objective evolutionary algorithm to search the space of GCC optimization flags and produces a Pareto front of settings that trade compilation time against execution time. The Pareto-optimal settings outperform the standard levels \code{-O1}, \code{-O2}, and \code{-O3}, and the paper argues that compiler developers should derive their optimization levels by such systematic exploration rather than by hand. The Pareto-front presentation of size against estimated cycles in this project's scatter plots and in the Phase-Ordering Explorer page follows the same multi-objective framing.

Nobre, Martins, and Cardoso \cite{nobre2016} proposed a graph-based approach to pass selection and ordering for LLVM. A directed graph over passes is built from sequences known to perform well, with edge weights reflecting how often one pass follows another; new sequences are then sampled as walks on the graph. The method finds sequences that improve on \code{-O3} for a range of kernels while exploring only a tiny fraction of the space. The idea that good orderings share local structure, that certain passes tend to follow certain others, is borne out by the structural rules found in this project, where every one of the best orderings shares the same relative positions of the four passes that matter.

\section*{Machine-learning approaches}

Agakov, Bonilla, Cavazos, Franke, Fursin, O'Boyle, Thomson, Toussaint, and Williams \cite{agakov2006} were among the first to use machine learning to speed up iterative compilation. Rather than predicting the best sequence directly, they learned a probability distribution over good transformation sequences from a training set of programs and used it to bias the random search on new programs, so that the search concentrated on the promising region of the space. On embedded targets the focused search reached the performance of an unfocused search an order of magnitude faster. The work established the pattern, followed by most later machine-learning approaches, of using static program features to transfer knowledge between programs.

Fursin and colleagues \cite{fursin2011} built this idea into a production compiler with Milepost GCC. The compiler extracts a vector of static features from each function, and a model trained on a corpus of programs predicts which optimization flags to enable. The system was integrated into a mainstream GCC release and demonstrated that learned optimization decisions could be deployed in practice, not only in research prototypes. The features it used, such as counts of instruction kinds, basic blocks, and edges, are the kind of features that this project's metric collector computes as an instruction-category breakdown.

Kulkarni and Cavazos \cite{kulkarnicavazos2012} reformulated phase ordering as a sequential decision problem inside a just-in-time compiler. A neural network, evolved with the NEAT neuroevolution algorithm, examines the current state of the method being compiled and decides which optimization to apply next, one step at a time, until it decides to stop. Because the decision is made from the current code rather than from the original program, the model adapts the sequence as the code changes. Applied to the Jikes research virtual machine, it produced consistent speedups over the default sequence. The stage-by-stage view in this project's Optimization Lab, where the effect of each pass on the current code is displayed before the next pass is chosen, is the interactive, human-driven counterpart of this formulation.

Ashouri, Bignoli, Palermo, Silvano, Kulkarni, and Cavazos \cite{ashouri2017} proposed MiCOMP, which mitigates the phase-ordering problem in LLVM by clustering the optimizer's passes into a small number of sub-sequences and then learning to predict a good ordering of the sub-sequences from program features. The reduction from individual passes to clusters shrinks the search space enough for prediction to be effective, and the predicted orderings outperformed \code{-O3} on a range of benchmarks. The clustering insight, that passes which always travel together can be treated as a unit, is echoed in this project's finding that certain passes never affect the outcome relative to each other and can be collapsed without loss.

Ashouri, Killian, Cavazos, Palermo, and Silvano \cite{ashouri2018} surveyed the field of compiler autotuning with machine learning. The survey classifies the literature by the problem addressed, selection of optimizations versus their ordering; by the learning technique, supervised, unsupervised, or reinforcement learning; by the program representation used as input; and by the target, and identifies open problems including the lack of standard benchmarks and the difficulty of reproducing results across compiler versions. The survey was used in this project as a map of the field and confirmed that small, transparent infrastructures for controlled experimentation are rare.

Haj-Ali, Huang, Moses, Xiang, Stoica, Asanovi\'{c}, and Wawrzynek \cite{hajali2020} introduced AutoPhase, which applies deep reinforcement learning to phase ordering for high-level synthesis, where LLVM passes are applied to a program before it is compiled to hardware and the objective is circuit performance. The agent observes features of the current program and selects the next pass; a random-forest analysis identifies which features matter. AutoPhase reported large improvements over \code{-O3} for the HLS benchmarks studied. The paper is notable for demonstrating that the phase-ordering problem is as relevant to hardware generation as to software, and for showing that reinforcement learning can discover orderings that domain experts would not have chosen.

Jain, Andaluri, VenkataKeerthy, and Upadrasta \cite{jain2022} proposed POSET-RL, which targets the two objectives that matter most for embedded systems, code size and execution time, simultaneously. The method derives a dependence graph from the passes in LLVM's \code{-Oz} size-optimizing pipeline, uses it to constrain the sequences a reinforcement-learning agent may propose, represents programs with learned IR embeddings, and rewards the agent for improving size without degrading time. The reported results improved on \code{-Oz} in code size while maintaining or improving execution time on standard benchmark suites. POSET-RL was the primary inspiration for this project's choice of metrics: static size and a proxy for time are measured for every ordering, and the trade-off between them is the central object of study.

Cummins, Wasti, Guo, Cui, Ansel, Gomez, Jain, Liu, Teytaud, Steiner, Tian, and Leather \cite{cummins2022} released CompilerGym, a set of reinforcement-learning environments for compiler optimization problems built in the style of OpenAI Gym. The LLVM environment exposes phase ordering as an environment whose actions are passes and whose rewards are changes in code size or runtime, together with standard benchmark datasets and baselines. The project made phase-ordering research reproducible and accessible to machine-learning researchers without compiler expertise. This project's web API, which exposes compile, optimize, step, and benchmark endpoints over HTTP, is a much smaller but philosophically similar attempt to make a phase-ordering environment available to others.

VenkataKeerthy, Aggarwal, Jain, Desarkar, Upadrasta, and Srikant \cite{venkatakeerthy2020} developed IR2Vec, a method for embedding LLVM intermediate representation into continuous vectors that capture the semantics of the code. The embeddings are computed from the IR itself, without executing the program, and are used as input features by POSET-RL and other learning-based phase-ordering systems. The work illustrates how much effort the machine-learning approaches must invest in representing programs before any learning can begin, and by contrast how simple the representation problem becomes in an infrastructure such as this project's, where the IR has 35 opcodes and the entire program is a short list of instructions.

Cummins, Seeker, Grubisic, Elhoushi, Liang, Roziere, Gehring, Gloeckle, Hazelwood, Synnaeve, and Leather \cite{cummins2023} trained a large language model to optimize LLVM assembly. Given an unoptimized function, the model emits the list of passes that should be applied to minimize code size, and as an auxiliary task it predicts the resulting instruction count and the optimized code itself. The model improved on \code{-Oz} by a few percent over a large test set without executing any compilation during inference. The work represents the current frontier of learning-based phase ordering and highlights that even the most capable models are being evaluated against exactly the metric this project uses, static instruction count.

Han and Kim \cite{han2024} observed that many passes in a sequence are dormant, changing nothing, and that a search algorithm which can predict dormancy can effectively explore an exponentially larger space of sequences for the same cost, since every dormant pass can be inserted or removed without evaluation. They learned to predict dormancy from program features and integrated the predictor into a phase-ordering search. This is a modern, learned version of the dormant-phase detection of Kulkarni et al.\ \cite{kulkarni2004}, and it is precisely the phenomenon that produces the large equivalence classes in this project's results.

Liu and colleagues \cite{liu2024} proposed to model the dependence between passes explicitly and to use the dependence model to prune the ordering search, so that passes known to be independent are not permuted and passes known to depend on one another are kept in the beneficial order. Their approach reduced search cost while preserving optimization quality. Chen, Chen, Yao, Yu, Xu, and Ma \cite{chen2025} identified short chains of passes that reliably boost the effectiveness of the passes that follow them, and built phase-ordering sequences around these booster chains. Both papers, published after this project began, articulate the same structural view that this project's results support empirically: the ordering space is governed by a small number of dependence relationships among a subset of the passes, and once those are respected the remaining freedom is inconsequential.

\section*{Infrastructure, targets, and cost models}

Lattner and Adve \cite{lattner2004} described the LLVM compiler infrastructure, whose modular pass framework and typed, static-single-assignment intermediate representation have made it the platform for almost all recent phase-ordering research. LLVM's pass manager accepts an explicit list of passes and runs them in the given order, which is what makes the phase-ordering experiments of \cite{ashouri2017,hajali2020,jain2022,cummins2023} possible. The pass manager in this project is modelled on the same idea at a much smaller scale: a registry of named passes, a validated ordered list, and a loop that threads the IR through each pass in turn. LLVM was not used directly because its scale, with over a hundred passes and an IR that a student cannot read at a glance, defeats the purpose of a transparent teaching and research tool.

Waterman and Asanovi\'{c} \cite{waterman2019} specify the RISC-V unprivileged instruction set architecture. RISC-V is an open, royalty-free ISA whose base integer instruction set, RV32I, together with the multiply and divide extension, M, defines the RV32IM target chosen for this project's back end. The specification's small, regular instruction set makes it possible to write a correct code generator and a validating simulator in a few hundred lines each, and its rapidly growing adoption in low-cost microcontrollers makes it a realistic embedded target. The specification also documents the behaviour of division by zero, which the back end must guard explicitly to match the interpreter's semantics.

The ARM Cortex-M3 Technical Reference Manual \cite{arm2010} provided the instruction timing information from which this project's cost model was derived. On Cortex-M class cores, most arithmetic, logical, and comparison instructions execute in a single cycle; multiplication takes one to several cycles depending on the core; hardware division takes between two and twelve cycles depending on the operands; taken branches incur a pipeline refill penalty of several cycles; and loads and stores take two cycles. The cost model assigns weights in this spirit, one cycle for simple operations, three for multiplication, twelve for division and modulo, two for branches and memory accesses, and larger fixed costs for calls, returns, and output. The model is an estimate, not a measurement, and Chapter 5 discusses its limitations.

Georgiou, Blackmore, Xavier-de-Souza, and Eder \cite{georgiou2018} studied the standard optimization levels of GCC and LLVM on Cortex-M embedded targets from the point of view of both performance and energy consumption. They found that the highest optimization level is frequently not the best choice on these targets, that lower levels sometimes produce faster and more energy-efficient code, and that the reasons lie in pass interactions and target-specific effects that the default pipelines do not account for. They proposed a lightweight approach that exploits the existing levels rather than searching the full flag space. Their work supports the premise of this project that the phase-ordering problem is of practical importance on embedded targets and that the default pipeline is not the last word.

\section*{Benchmark suites and evaluation methodology}

Guthaus, Ringenberg, Ernst, Austin, Mudge, and Brown \cite{guthaus2001} introduced MiBench, a free benchmark suite designed to be representative of commercial embedded workloads. The suite is organized into six categories, automotive and industrial control, consumer devices, office automation, networking, security, and telecommunications, and its programs are small enough to run on a simulator yet realistic enough that their characteristics differ measurably from the desktop-oriented SPEC suites. Eight of this project's fifteen benchmarks are adapted from MiBench kernels: bit counting from the automotive category, integer square root and basic arithmetic kernels from basicmath, and modular exponentiation and hash mixing from the security category. The adaptations preserve the control structure and arithmetic character of the originals while restricting them to the integer subset of C that the compiler supports.

Pouchet and Yuki \cite{pouchet2016} maintain PolyBench/C, a collection of numerical kernels from linear algebra, stencil computations, data mining, and medley categories, written in a restricted style that makes them amenable to polyhedral analysis. Each kernel is a small set of nested loops over arrays with affine index expressions. Seven of this project's benchmarks are adapted from PolyBench: the one- and two-dimensional Jacobi stencils, general matrix multiplication (GEMM), two chained matrix multiplications (2MM), matrix-transpose-vector product (ATAX), the BiCG sub-kernel, and the scalar-vector-matrix product GESUMMV. These kernels were added in Phase II to test whether the phase-ordering effects observed on scalar MiBench kernels carry over to array-based numerical code with repeated index arithmetic, which they do.

Fleming and Wallace \cite{fleming1986} showed, in a widely cited note, that the arithmetic mean is the wrong way to summarize normalized benchmark results, because the ranking it produces depends on which system is chosen as the normalization base. The geometric mean is the only mean that is consistent under normalization, and it is therefore the correct aggregate for ratios such as ``optimized size divided by baseline size''. All cross-program summaries in this project are geometric means of per-program ratios, computed as the exponential of the mean logarithm, and every result table that reports an average across programs follows this rule.

\section*{Educational compiler tools}

Godbolt \cite{godbolt2019} described Compiler Explorer, the widely used web tool that shows, for a program typed into the browser, the assembly produced by dozens of compilers at any optimization level, with source lines colour-linked to assembly lines. The article walks through the optimizations a modern compiler performs and shows how the tool reveals them. Compiler Explorer inspired the side-by-side, hover-linked presentation of intermediate code and RISC-V assembly in this project's Assembly page. It differs from this project in that it treats the compiler as a black box with a fixed pipeline; the user can choose \code{-O2} but cannot reorder the passes within it or see the intermediate stages.

Koitz-Hristov, Mandl, and Wotawa \cite{koitz2025} presented VisOpt, an educational tool that visualizes compiler optimizations for computer science students, showing the effect of individual transformations on program representations to support teaching of compiler construction. Their evaluation in a classroom setting found that visualization improved students' understanding of what the optimizations do. This project shares VisOpt's goal and extends it in one specific direction: it visualizes not only individual optimizations but their interaction, by letting the learner change the order and observe the consequence at each stage.

Nystrom \cite{nystrom2021} and Sandler \cite{sandler2024} are recent book-length tutorials on building interpreters and C compilers from scratch. Both emphasize hand-written lexers and recursive-descent parsers, incremental development with a test suite, and a three-address or similar linear intermediate form, and Sandler's book in particular covers the implementation of constant folding, copy propagation, unreachable-code elimination, and dead-store elimination as separate passes over a three-address code, together with a discussion of the order in which to run them. These works guided the engineering style of this project's front end and confirmed that a compiler of this scope can be implemented, tested, and understood by a small student team.

\section*{Summary and research gap}

Table~\ref{tab:lit-summary} summarizes the surveyed work by approach, the compiler on which it was carried out, the scale of the pass set, and the objective measured.

\begin{table}[!tbp]
\centering\footnotesize
\caption{Summary of the surveyed work on phase ordering and compiler autotuning}
\label{tab:lit-summary}
\begin{tabular}{L{3.1cm} L{2.7cm} L{2.3cm} L{1.7cm} L{2.7cm}}
\toprule
\textbf{Work} & \textbf{Approach} & \textbf{Compiler} & \textbf{Passes} & \textbf{Objective}\\
\midrule
Whitfield \& Soffa \cite{whitfield1997} & Formal enabling / disabling relations & Genesis & few & interaction analysis\\
Cooper et al.\ \cite{cooper1999} & Genetic algorithm & research (MSCP) & 10 & code size\\
Almagor et al.\ \cite{almagor2004} & Exhaustive sub-space, hill climbing & research & 5 (len.\ 10) & speed, size\\
Kulkarni et al.\ \cite{kulkarni2004} & GA with dormancy and hashing & VPO / VISTA & 15 & size, speed\\
Kulkarni et al.\ \cite{kulkarni2006,kulkarni2009} & Exhaustive instance enumeration & VPO & 15 & size, dyn.\ count\\
Triantafyllis et al.\ \cite{triantafyllis2003} & Configuration exploration, static estimator & Itanium research & configurations & speed\\
Hoste \& Eeckhout \cite{hoste2008} & Multi-objective evolutionary & GCC & flags & Pareto speed / compile time\\
Agakov et al.\ \cite{agakov2006}, Fursin et al.\ \cite{fursin2011} & ML-focused search, learned flag prediction & GCC & flags & speed\\
Kulkarni \& Cavazos \cite{kulkarnicavazos2012} & Neuroevolution, per-step decision & Jikes RVM & tens & speed\\
Ashouri et al.\ \cite{ashouri2017} & Clustered sub-sequences, supervised ML & LLVM & $\sim$100 & speed\\
Haj-Ali et al.\ \cite{hajali2020} & Deep RL & LLVM / HLS & 45 & circuit perf.\\
Jain et al.\ \cite{jain2022} & Deep RL with dependence graph & LLVM & $\sim$30 & size and time\\
Cummins et al.\ \cite{cummins2022,cummins2023} & RL environments; LLM pass prediction & LLVM & $\sim$120 & size\\
Han \& Kim \cite{han2024}, Liu et al.\ \cite{liu2024}, Chen et al.\ \cite{chen2025} & Dormancy, dependence, booster chains & LLVM & $\sim$100 & size, speed\\
Georgiou et al.\ \cite{georgiou2018} & Standard-level analysis on Cortex-M & GCC, LLVM & levels & speed, energy\\
Godbolt \cite{godbolt2019}, Koitz-Hristov et al.\ \cite{koitz2025} & Visualization for learners & GCC, LLVM, others & fixed & education\\
\textbf{This project} & \textbf{Exhaustive 721-ordering enumeration, equivalence-class analysis, interactive tool} & \textbf{custom, transparent} & \textbf{6} & \textbf{size, est.\ cycles, dyn.\ count}\\
\bottomrule
\end{tabular}
\end{table}

Three observations follow from the survey. First, almost all quantitative phase-ordering research has been conducted on production compilers, chiefly LLVM, GCC, and VPO. This has the advantage of realism but the disadvantage of opacity: with a hundred passes and a complex IR, the researcher measures the aggregate effect of ordering but cannot easily attribute it to specific pairwise interactions, and a student cannot follow what happened. Second, the search-based and learning-based methods are motivated by the impossibility of exhaustive enumeration at production scale, and they therefore report how much better than the default a found sequence is, but rarely how far from the true optimum it remains; only the exhaustive studies of Kulkarni and colleagues can answer that question, and they did so for a fixed set of fifteen phases in a single research compiler. Third, educational tools show the effect of a fixed pipeline or of single optimizations, but none allows a learner to reorder passes and observe the consequences stage by stage.

The gap that this project addresses lies at the intersection of these three observations. It builds a compiler infrastructure small enough that every pass and every intermediate result can be read, in which the ordering space is small enough to enumerate completely, so that the true optimum and the structure of the space are known exactly rather than estimated; it evaluates on kernels adapted from the embedded and numerical benchmark suites used by the wider community, with the metrics, size, time proxy, and dynamic count, that the community uses; it lowers the results to a real instruction set to check that improvements in the IR survive code generation; and it wraps the whole in an interactive tool that makes the phase-ordering problem visible to learners. None of the surveyed works combines these properties.
